\section{The finite rational certificate}
\label{app:certificate}
This appendix specifies the seven rows and all arithmetic needed to verify
\eqref{eq:3certificate} and \eqref{eq:4certificate}. Starting from the integer
rows in Section~\ref{app:rows}, evaluate the rational three-body traces in
Section~\ref{app:three-arithmetic} and the rational four-body matrices in
Section~\ref{app:four-arithmetic}. The margin and coefficient-sign tables in
those sections record the resulting checks. The formulas below determine
every value using only integers and rational numbers.

\subsection{Row data}
\label{app:rows}
Let $P=10^7$. The first two numbers of each row are $t$ and $\ell=P\lambda$. An entry $pj:a$ specifies the integer $a_{pj}$, with $p$ and $j$ parsed as separate single-digit labels; omitted entries are zero. Continuation lines belong to the preceding row. Define
\begin{equation}\label{eq:alphadata}
 \alpha_{pj}=\frac{a_{pj}}P
       \sqrt{\frac{2^p\binom{p+j}{p}}{2^t\binom{p+j}{t}}}
 =\frac{a_{pj}}P\sqrt{\frac{2^{p-t}t!(p+j-t)!}{p!j!}}.
\end{equation}
Every nonzero entry has $0\le p+j-t\le8$.
\par\noindent\begin{minipage}{\linewidth}
\lstinputlisting{../verification/rows.txt}
\end{minipage}\par
\smallskip
\noindent In particular $P^2\eta=\sum_\rho\ell_\rho^2=93527408868499$.

\subsection{Three-body arithmetic}
\label{app:three-arithmetic}
For $c\in\{1,2\}$, an integer $T\ge0$, and integers $z,p$, define
\begin{equation}\label{eq:U}
 \mathsf U(c,z,T,p)=
 \sum_{h=\max(0,p+z-T)}^{\min(p,z)}
 (-1)^{z-h}\binom zh\binom{T-z}{p-h}c^{p-h},
\end{equation}
and set it to zero if $z$ or $p$ lies outside $[0,T]$. Expanding \eqref{eq:coupling} gives
\[
 s^{(1/\sqrt3)}_{z,T,p}=\mathsf U(2,z,T,p)
 \sqrt{\frac{2^z\binom Tz}{3^T2^p\binom Tp}}.
\]
Substitution in \eqref{eq:ez} and \eqref{eq:alphadata} cancels all radicals:
\begin{equation}\label{eq:erational}
 P^2e_z=\sum_\rho\sum_{T=\max(z,t_\rho)}^{15}
 \frac{2^z\binom Tz}{3^T2^{t_\rho}\binom T{t_\rho}}
 X_{\rho,T}\left(X_{\rho,T}+2\ell_\rho\mathsf U(2,z,T,t_\rho)\right),
\end{equation}
where $X_{\rho,T}=\sum_{p+j=T}a^\rho_{pj}\mathsf U(2,z,T,p)$, with value zero on empty levels. For each $z\in\{2,4,5,\ldots,15\}$, evaluate \eqref{eq:erational} and form the exact rational margin
$m_z=\frac32(3z-1)(-1/2)^z-e_z$. Write $10^6m_z=n_z/d_z$ with integers $n_z,d_z$ and $d_z>0$. Integer division gives the floors $\lfloor n_z/d_z\rfloor$ listed below:
\begin{center}
\begin{tabular}{rrrrrrrrrrrrrr}
\toprule
$z$&2&4&5&6&7&8&9&10&11&12&13&14&15\\
\midrule
floor&1298313&211317&249&250&249&250&250&250&249&250&250&6470&250\\
\bottomrule
\end{tabular}
\end{center}
All are positive, proving \eqref{eq:3certificate}.

\subsection{Four-body arithmetic}
\label{app:four-arithmetic}
Fix $1\le D\le23$; all copy indices below lie in $\mathcal R_D$. Put
\[
 u_r=\frac1{r!(D-r)!},\qquad
 (E_D)_{rs}=\sqrt{u_ru_s}\,Y_{rs},\qquad
 (G_D)_{rs}=\sqrt{u_ru_s}\,Z_{rs}.
\]
The matrices $Y,Z$ are rational. Define
\[
 \mathsf V_r(D,T;p,j,k)=
 \mathsf U(1,r,j+k,j)\mathsf U(1,D-r,T-r,p),
\]
with value zero for $r>j+k$. Two applications of \eqref{eq:coupling} give
\begin{equation}\label{eq:Srational}
 S_r(D,T;p,j,k)=\mathsf V_r(D,T;p,j,k)
 \sqrt{2^{1-(j+k)-T+r}\frac{j!k!p!\,u_r}{(T-D)!}}.
\end{equation}
Therefore
\begin{equation}\label{eq:Yrational}
 \begin{split}
 P^2Y_{rs}=-\sum_\rho\sum_{\substack{p,j;q,l\\T\ge D}}
 &a^\rho_{pj}a^\rho_{ql}\,
 \mathsf V_r(D,T;p,j,k)\mathsf V_s(D,T;q,l,i)\\
 &\times2^{1-2T+p+q-t_\rho+(r+s)/2}
       \frac{i!k!t_\rho!}{(T-D)!},
 \end{split}
\end{equation}
where $i=p+j-t_\rho$, $k=q+l-t_\rho$, and $T=p+j+k$. Both entry sums are ordered. To verify the cancellation, multiply the two factors in \eqref{eq:Srational} by the two factors \eqref{eq:alphadata}; the factorial product under the combined square root is
$(t_\rho!i!k!)^2u_ru_s/((T-D)!)^2$.
Using $j+k=T-p$ and $l+i=T-q$, its power of two has exponent
$2+2p+2q-2t_\rho-4T+r+s$. Taking the positive square root and
removing $\sqrt{u_ru_s}$ gives exactly the displayed factor.

For an increasing quadruple $A$ of nonnegative indices with sum $D+1$, put $f(A)=\prod_{b\in A}b!$. If $(x,y,j,k)$ lists its elements, let $\sigma(x,y,j,k)$ be the sign of sorting this list into increasing order. In each summand below put $p=x+y-1$, and define
\begin{equation}\label{eq:Lrational}
 L_r(A)=2^{(1-2D+r)/2}
 \sum_{\substack{x<y\text{ in }A\\\{j,k\}=A\setminus\{x,y\},\ j<k}}
 \sigma(x,y,j,k)(x-y)\,p!j!k!\,
 \mathsf V_r(D,D;p,j,k).
\end{equation}
Combining \eqref{eq:pairstar} and \eqref{eq:Srational}, the coefficient of $w_{Dr}^*$ on the Slater wedge $A$ is $\sqrt{u_r/f(A)}L_r(A)$. Thus
\begin{equation}\label{eq:Zrational}
 Z_{rs}=\sum_{\substack{A\text{ increasing quadruple}\\\sum_{b\in A}b=D+1}}
                   \frac{L_r(A)L_s(A)}{f(A)}.
\end{equation}
All exponents of two in these formulas are integers, since $r,s$ are odd.

Let $D_u=\diag(u_r)$ and $D_{\sqrt u}=\diag(\sqrt{u_r})$. The matrix in \eqref{eq:4certificate} equals $D_{\sqrt u}MD_{\sqrt u}$, where
\begin{equation}\label{eq:M}
 M=ZBZ,\qquad
 B=\diag\bigl((2\mathbf1_{r=1}+\eps)u_r\bigr)-D_uYD_u.
\end{equation}
Thus it suffices to check $M\ge0$ using rational arithmetic. This
congruence uses only the positive diagonal $D_{\sqrt u}$; it does not
invert $G_D$ or presume that its columns are independent.

To check positivity for each $D=1,\ldots,23$, form $Y$ and $Z$ from
\eqref{eq:Yrational} and \eqref{eq:Zrational}, then form $B$ and $M$ from
\eqref{eq:M}, using $\eps=3/10^6$. All sums range over the printed rows and
the finite index sets specified above. Compute with exact fractions
throughout. If $d=|\mathcal R_D|$, initialize $J=I_d$. For
$b=1,\ldots,d$, compute in order
\begin{equation}\label{eq:recurrence}
 X=MJ,\qquad c_b=\frac{\tr X}{b},\qquad J=c_bI_d-X.
\end{equation}
This is the Leverrier--Faddeev recurrence applied to $-M$;
see \cite{HouConference}. It gives
$\det(xI+M)=x^d+c_1x^{d-1}+\cdots+c_d$. For example, this recurrence follows by multiplying the polynomial adjugate of $xI+M$ by $xI+M$ and comparing coefficients, together with the derivative-of-determinant trace identity. Exact evaluation of \eqref{eq:U}, \eqref{eq:Yrational}, \eqref{eq:Lrational}, \eqref{eq:Zrational}, and \eqref{eq:recurrence} gives:
\begin{center}
\begin{tabular}{rl@{\qquad}rl@{\qquad}rl}
\toprule
$D$&signs&$D$&signs&$D$&signs\\
\midrule
1&\texttt{0}&9&\texttt{++000}&17&\texttt{++++++000}\\
2&\texttt{0}&10&\texttt{+0000}&18&\texttt{+++++0000}\\
3&\texttt{00}&11&\texttt{+++000}&19&\texttt{+++++++000}\\
4&\texttt{00}&12&\texttt{++0000}&20&\texttt{++++++0000}\\
5&\texttt{+00}&13&\texttt{++++000}&21&\texttt{++++++++000}\\
6&\texttt{000}&14&\texttt{+++0000}&22&\texttt{+++++++0000}\\
7&\texttt{+000}&15&\texttt{+++++000}&23&\texttt{+++++++++000}\\
8&\texttt{+000}&16&\texttt{++++0000}&&\\
\bottomrule
\end{tabular}
\end{center}
The signs list the coefficients after the leading $1$, in descending degree.
For each computed coefficient, choose a positive denominator and inspect
its integer numerator: \texttt{+} denotes a positive numerator and
\texttt{0} denotes zero. We use the characteristic-coefficient positivity test for symmetric
matrices~\cite[Theorem~5]{PeyrlParrilo}. Every coefficient is nonnegative,
so $\det(xI+M)>0$ for $x>0$. Since $M$ is real symmetric, a negative eigenvalue would give a positive root of this polynomial. Therefore $M\ge0$, proving \eqref{eq:4certificate}.

\begin{remark}[Nature of the certificate]
The data above certify fixed spin-matrix inequalities. The positive comparison operator is a sum of squares by construction. No many-particle finite-size spectrum, numerical extrapolation, or assumed plasma property enters the argument. The exact equalities and inequalities in this appendix can be verified by a finite calculation using only integers and rational numbers.
\end{remark}
